\chapter{The Periodic Table of Elements: A Comprehensive SDT Derivation}


%----------------------------------------------------------------------
\section{Introduction: From Abstract Rules to a Geometric Catalogue}
%----------------------------------------------------------------------

The periodic table organises the fundamental building blocks of matter into a coherent, predictive pattern.  SDT contends that this pattern is a direct, visual representation of the \textbf{stable, geometric solutions for packing resonant displacement vortices.}

This chapter serves as a comprehensive catalogue and derivation of element properties from first principles.


%----------------------------------------------------------------------
\section{Methodology: The SDT Signature}
%----------------------------------------------------------------------

For each element, four parameters define its ``SDT Signature'':

\begin{enumerate}
  \item \textbf{Input:} atomic number $Z$ and mass number $A$ (most stable isotope).
  \item \textbf{$R_{\text{nucleus}}$:} Nuclear radius from $R \approx 1.2 \times A^{1/3}$\,fm.
  \item \textbf{$k_{\text{nucleus}}$:} Nuclear k-factor, $k_{\text{nucleus}} \approx k_p / A^{1/3}$.
  \item \textbf{$k_{\text{atom}}$:} Atomic k-factor from first ionisation energy: $k = c/\sqrt{2 E_{I1}/m_e}$.
  \item \textbf{$\chi_{\text{SDT}}$:} SDT electronegativity $\propto P_{\text{surface}} \times (R_{\text{nucleus}}/r_{\text{covalent}})^2$.
\end{enumerate}


%----------------------------------------------------------------------
\section{Selected Element Derivations}
%----------------------------------------------------------------------

\begin{center}
\begin{tabular}{rlllrrrrl}
\hline
$Z$ & \textbf{El.} & \textbf{Config.} & $R_{\text{nuc}}$ & $k_{\text{nuc}}$ & $E_{I1}$ & $k_{\text{atom}}$ & $\chi$ & \textbf{Character} \\
 & & & (fm) & & (eV) & & & \\
\hline
 1 & H  & $1s^1$              & 1.20 & 0.546 & 13.60 & 137.0 & 2.20 & Reactive nonmetal \\
 2 & He & $1s^2$              & 1.90 & 0.344 & 24.59 & 101.9 & --- & Sealed shell \\
\hline
 3 & Li & [He]$2s^1$          & 2.30 & 0.285 &  5.39 & 217.9 & 0.98 & Alkali metal \\
11 & Na & [Ne]$3s^1$          & 3.29 & 0.192 &  5.14 & 223.1 & 0.93 & Alkali metal \\
19 & K  & [Ar]$4s^1$          & 3.86 & 0.163 &  4.34 & 242.7 & 0.82 & Alkali metal \\
\hline
 9 & F  & [He]$2s^22p^5$      & 3.12 & 0.205 & 17.42 & 121.1 & 3.98 & Halogen \\
17 & Cl & [Ne]$3s^23p^5$      & 3.75 & 0.167 & 12.97 & 140.4 & 3.16 & Halogen \\
35 & Br & [Ar]$4s^23d^{10}4p^5$ & 4.90 & 0.136 & 11.81 & 147.1 & 2.96 & Halogen \\
\hline
26 & Fe & [Ar]$4s^23d^6$      & 4.49 & 0.147 &  7.90 & 179.9 & 1.83 & Transition metal \\
47 & Ag & [Kr]$5s^14d^{10}$   & 5.28 & 0.125 &  7.58 & 183.7 & 1.93 & Transition metal \\
79 & Au & [Xe]$6s^14f^{14}5d^{10}$ & 5.82 & 0.109 &  9.23 & 166.4 & 2.54 & Noble metal \\
\hline
92 & U  & [Rn]$7s^25f^36d^1$  & 6.18 & 0.103 &  6.19 & 202.9 & 1.38 & Actinide \\
\hline
\end{tabular}
\end{center}


%----------------------------------------------------------------------
\section{Analysis of Periodic Trends}
%----------------------------------------------------------------------

\subsection{Atomic Radius}

Defined by the outermost stable resonant shell ($r_n$).  Correctly increases down a group (increasing $n$) and decreases across a period (increasing $Z_{\text{eff}}$).

\subsection{Ionisation Energy and $k_{\text{atom}}$}

Inversely related.  High $E_{I1}$ means tight binding, high velocity, low $k_{\text{atom}}$.  $k_{\text{atom}}$ is highest for alkali metals (low $E_{I1}$) and lowest for noble gases and halogens (high $E_{I1}$).

\subsection{Electronegativity ($\chi_{\text{SDT}}$)}

A direct measure of the propagated pressure field at the covalent boundary.  Highest for small atoms with high nuclear charge (F); lowest for large atoms with low effective valence charge (K).

\subsection{Metallic Character}

A direct consequence of \textbf{high $k_{\text{atom}}$}.  Atoms with high k-factors have ``slow,'' loosely bound valence electrons that easily detach to form a \textbf{delocalised sea of electron vortices}---the metallic bond.  Metallic character decreases left-to-right as k decreases.


%----------------------------------------------------------------------
\section{The d and f Blocks: Complex Resonance Geometry}
%----------------------------------------------------------------------

\textbf{The d-block dip:} the 4s shell fills before 3d because a simple, spherically symmetric 4s vortex is a more stable, lower-energy configuration than a complex 3d vortex when the nuclear charge is insufficient to support the latter.

Only at $Z = 21$ (Sc) does the $n = 3$ shell become ``tight'' enough for a stable 3d resonance.

This is a competition between stable geometric solutions, not an arbitrary rule.  The atom settles into the configuration of minimum total integrated pressure.


%----------------------------------------------------------------------
\section{Conclusion}
%----------------------------------------------------------------------

The periodic table is a \textbf{catalogue of the stable geometric and resonant solutions for $N$-body vortex systems}.  The trends in atomic radius, ionisation energy, electronegativity, and metallic character all emerge as consequences of the changing size and pressure of the nuclear vortex and the geometric packing rules of the electron vortices.

The seemingly complex structure of the table is an elegant map of the solutions to a single physical problem: \textbf{finding the minimum-energy configuration for a set of displacement vortices in a pressurised medium.}
